% !TEX program = pdflatex
\documentclass[11pt,a4paper]{article}

% =========================================================================
% ESPECIFICACIÓN DE REQUERIMIENTOS
% Proyecto: Horus Sentinel
% Gestión de Proyectos TICS (EST315) -- Semestre 2026-II -- Team Sobrevivientes
% =========================================================================

% --- PAQUETES DE CONFIGURACIÓN ---
\usepackage[utf8]{inputenc}
\usepackage[T1]{fontenc}
\usepackage[spanish,es-tabla]{babel}
% Compatibilidad con babel-spanish reciente (flechas de TikZ y signo de porcentaje)
\AtBeginDocument{\shorthandoff{<>}\spanishplainpercent}
\usepackage{geometry}
\geometry{
    top=3.6cm,
    bottom=2.3cm,
    left=2.2cm,
    right=2.2cm,
    % PARTE MODIFICABLE: headheight debe ser mayor que la altura del logo más alto (2.56cm = 73pt).
    headheight=80pt,
    headsep=15pt,
    footskip=20pt
}

\usepackage{graphicx}
\usepackage{fancyhdr}
\usepackage[table]{xcolor}
\usepackage{titlesec}
\usepackage{booktabs}
\usepackage{tabularx}
\usepackage{array}
\usepackage{amsmath}
\usepackage{enumitem}
\usepackage{microtype}
\usepackage{hyperref}

% Logos: funciona a 2 o 3 niveles de profundidad dentro de docs/latex/
\graphicspath{{../../assets/img/}{../../../assets/img/}}

% --- PALETA INSTITUCIONAL FIEI / UNA PUNO ---
\definecolor{AzulUNA}{RGB}{10, 45, 90}
\definecolor{DoradoFIEI}{RGB}{180, 130, 20}
\definecolor{GrisOscuro}{RGB}{50, 50, 50}
\definecolor{AzulClaro}{RGB}{228, 236, 246}

% --- CONFIGURACIÓN DE LISTAS ---
\setlist{topsep=2pt, itemsep=2pt, parsep=1pt}

% --- COLUMNAS DE TABLA ---
\newcolumntype{L}{>{\raggedright\arraybackslash}X}
\newcolumntype{P}[1]{>{\raggedright\arraybackslash}p{#1}}
\newcolumntype{C}[1]{>{\centering\arraybackslash}p{#1}}
\renewcommand{\arraystretch}{1.25}

% --- FORMATO DE ENLACES ---
\hypersetup{
    colorlinks=true,
    linkcolor=AzulUNA,
    urlcolor=AzulUNA,
    citecolor=AzulUNA,
    pdftitle={Especificación de Requerimientos -- Horus Sentinel},
    pdfauthor={Team Sobrevivientes}
}

% --- FORMATO DE TÍTULOS ---
\titleformat{\section}{\color{AzulUNA}\normalfont\large\bfseries}{\thesection.}{0.7em}{}[\color{AzulUNA}\titlerule]
\titleformat{\subsection}{\color{AzulUNA}\normalfont\normalsize\bfseries}{\thesubsection.}{0.5em}{}
\titlespacing*{\section}{0pt}{9pt}{4pt}
\titlespacing*{\subsection}{0pt}{6pt}{2pt}

% --- ENCABEZADO Y PIE DE PÁGINA (fancyhdr) ---
\pagestyle{fancy}
\fancyhf{}

% ===== TAMAÑO DE LOS LOGOS (PARTE MODIFICABLE) =====
% logo_fiei.png  : height=1.45cm (lado izquierdo)
% logo_equipo.png: height=2.56cm (lado derecho; 25 % más grande porque la imagen tiene más borde blanco)
% Para cambiar el tamaño, modifique solo el valor height=... de cada \includegraphics.
% Si el logo más alto supera 2.6cm, aumente también headheight y top en la configuración de página.
\fancyhead[L]{%
    \includegraphics[height=1.45cm,keepaspectratio]{logo_fiei.png}%
}
\fancyhead[C]{%
    \scriptsize \textbf{\textcolor{AzulUNA}{UNIVERSIDAD NACIONAL DEL ALTIPLANO DE PUNO}}\\[2pt]
    \tiny \textcolor{GrisOscuro}{FACULTAD DE INGENIERÍA ESTADÍSTICA E INFORMÁTICA}\\[2pt]
    \scriptsize \textbf{\textcolor{AzulUNA}{Gestión de Proyectos TICS (EST315) \textbar\ Team Sobrevivientes}}\\[2pt]
    \footnotesize \textbf{\textcolor{DoradoFIEI}{ESPECIFICACIÓN DE REQUERIMIENTOS}}%
}
\fancyhead[R]{%
    \includegraphics[height=2.56cm,keepaspectratio]{logo_equipo.png}%
}

\renewcommand{\headrulewidth}{0pt}
\renewcommand{\headrule}{}

\fancyfoot[L]{\scriptsize \textcolor{GrisOscuro}{Horus Sentinel}}
\fancyfoot[C]{\small \thepage}
\fancyfoot[R]{\scriptsize \textcolor{GrisOscuro}{Puno, Perú -- Semestre 2026-II}}
\renewcommand{\footrulewidth}{0.4pt}

\begin{document}

\begin{center}
    {\Large \textbf{\textcolor{AzulUNA}{ESPECIFICACIÓN DE REQUERIMIENTOS}}}\\[3pt]
    {\normalsize \textbf{Horus Sentinel -- Sistema de mapeo predictivo de zonas de riesgo delictivo}}\\[4pt]
    \footnotesize \textbf{Fecha:} viernes 04/09/2026 (semana 2) \quad \textbar\quad \textbf{Versión:} 1.0 \quad \textbar\quad \textbf{Fase:} Inicio
\end{center}

\vspace{0.1cm}

\section{Datos Generales}
\noindent\begin{tabularx}{\textwidth}{@{}P{4.2cm} L@{}}
\toprule
\textbf{Proyecto:} & Horus Sentinel \\
\textbf{Responsable:} & Eliseo Tarqui Ajahuana -- Especialista Técnico / Arquitectura \\
\textbf{Zona piloto:} & Un distrito de la ciudad de Puno \\
\textbf{Producto de la Unidad I:} & Prototipo funcional con datos simulados \\
\bottomrule
\end{tabularx}

\section{Descripción General}
\begin{itemize}[leftmargin=*]
    \item \textbf{Problema:} no existe información clara y actualizada sobre qué zonas son más peligrosas y en qué horarios; los desplazamientos y el patrullaje se deciden sin datos.
    \item \textbf{Objetivo del sistema:} mostrar un mapa de calor de riesgo delictivo y sugerir la ruta más segura entre dos puntos de la zona piloto.
\end{itemize}

\section{Usuarios del Sistema}
{\small
\noindent\begin{tabularx}{\textwidth}{@{}P{3.4cm} L@{}}
\toprule
\rowcolor{AzulClaro}
\textbf{Usuario} & \textbf{Qué hace en el sistema} \\
\midrule
Ciudadano o estudiante & Consulta el mapa de riesgo y pide una ruta segura; en versiones posteriores, reporta incidentes. \\
Serenazgo o Policía & Consulta el mapa para priorizar el patrullaje. \\
Administrador & Carga y valida los datos de incidentes. \\
\bottomrule
\end{tabularx}
}

\section{Requerimientos Funcionales}
Prioridad: \textbf{D} = debe, \textbf{R} = debería, \textbf{P} = podría. Versión: \textbf{Prot.} = prototipo de la Unidad I con datos simulados; \textbf{MVP} = versión con datos reales.

\vspace{0.1cm}
{\small
\setlength{\tabcolsep}{4pt}
\noindent\begin{tabularx}{\textwidth}{@{}P{1.55cm} P{3.2cm} L C{0.6cm} C{0.9cm}@{}}
\toprule
\rowcolor{AzulClaro}
\textbf{ID} & \textbf{Requerimiento} & \textbf{Descripción y criterio de aceptación} & \textbf{Pri.} & \textbf{Ver.} \\
\midrule
RF-H01 & Conjunto de datos simulados & Generación de incidentes simulados en la zona piloto con tipo, coordenadas, fecha, hora y gravedad, siguiendo una distribución realista por zonas y horarios. & D & Prot. \\
RF-H02 & Mapa de calor de riesgo & Superficie de riesgo calculada con estimación de densidad por kernel (KDE) y mostrada sobre un mapa abierto de la ciudad. & D & Prot. \\
RF-H03 & Selección de origen y destino & El usuario marca en el mapa el punto de partida y el de llegada. & D & Prot. \\
RF-H04 & Ruta segura & Cálculo de la ruta sobre el grafo de calles con el algoritmo $A^*$ y costo $= \alpha \cdot \text{distancia} + \beta \cdot \text{riesgo}$; se muestra junto a la ruta más corta para compararlas. & D & Prot. \\
RF-H05 & Capas por horario & Mapa de riesgo filtrado por mañana, tarde, noche y madrugada. & R & MVP \\
RF-H06 & Reporte ciudadano & Formulario para reportar un incidente (tipo, lugar y hora) de forma anónima y moderada. & R & MVP \\
RF-H07 & Zonas anómalas & Alerta cuando una zona supera su nivel habitual de incidentes. & R & MVP \\
RF-H08 & Modelo espaciotemporal & Proceso de Poisson no homogéneo para estimar cómo un incidente eleva el riesgo cercano. & P & MVP \\
\bottomrule
\end{tabularx}
}

\section{Requerimientos No Funcionales}
{\small
\noindent\begin{tabularx}{\textwidth}{@{}P{1.85cm} P{2.8cm} L@{}}
\toprule
\rowcolor{AzulClaro}
\textbf{ID} & \textbf{Atributo} & \textbf{Requerimiento medible} \\
\midrule
RNF-H01 & Precisión espacial & Cada incidente se ubica con un margen de 15 metros. \\
RNF-H02 & Rendimiento & La ruta segura se calcula en menos de 2 segundos dentro de la zona piloto. \\
RNF-H03 & Privacidad & Solo se manejan incidentes agregados por zona y horario; no se identifican personas ni se usan cámaras. \\
RNF-H04 & Prevención del sesgo & El mapa indica la fuente y la fecha de los datos, y el modelo se audita para no sobrerrepresentar zonas más vigiladas. \\
RNF-H05 & Transparencia & El riesgo se presenta como una estimación, no como una certeza. \\
RNF-H06 & Costo & El prototipo usa mapas abiertos y servicios en capa gratuita. \\
\bottomrule
\end{tabularx}
}

\section{Modelo de Datos Principal}
{\small
\noindent\begin{tabularx}{\textwidth}{@{}P{3.0cm} L@{}}
\toprule
\rowcolor{AzulClaro}
\textbf{Entidad} & \textbf{Datos} \\
\midrule
Incidente & Tipo, coordenadas, fecha, hora, franja horaria, gravedad y fuente (simulado, reporte u oficial). \\
Zona & Identificador, límites y nivel de riesgo calculado por franja horaria. \\
Calle (grafo) & Nodos, tramos, longitud y riesgo asociado a cada tramo. \\
Ruta & Origen, destino, tramos recorridos, distancia total y riesgo total. \\
\bottomrule
\end{tabularx}
}

\section{Restricciones, Supuestos y Exclusiones}
\begin{itemize}[leftmargin=*]
    \item \textbf{Restricciones:} los datos georreferenciados de la Policía son reservados; el prototipo trabaja solo con datos simulados.
    \item \textbf{Supuestos:} existe un mapa abierto actualizado de las calles de la zona piloto.
    \item \textbf{Exclusiones:} no hay cobertura nacional, ni identificación de personas, ni vigilancia por cámaras en tiempo real.
\end{itemize}

\section{Firmas}
\noindent En señal de conformidad con la presente especificación, los integrantes del equipo firman a continuación:

\vspace{1.5cm}
\begin{center}
\begin{tabular}{c c c}
    \parbox{6.5cm}{\centering \rule{6.0cm}{0.6pt}\\[5pt] \textbf{Clyde Neil Paricahua Pari}\\[2pt] {\footnotesize Project Manager}}
    & \hspace{0.8cm} &
    \parbox{6.5cm}{\centering \rule{6.0cm}{0.6pt}\\[5pt] \textbf{Ruth Karina Apaza Solis}\\[2pt] {\footnotesize Analista de Planificación}}
    \\[2.0cm]
    \parbox{6.5cm}{\centering \rule{6.0cm}{0.6pt}\\[5pt] \textbf{Eliseo Tarqui Ajahuana}\\[2pt] {\footnotesize Especialista Técnico / Arquitectura}}
    & \hspace{0.8cm} &
    \parbox{6.5cm}{\centering \rule{6.0cm}{0.6pt}\\[5pt] \textbf{Viviana Mary Quisbert Quispe}\\[2pt] {\footnotesize Gestora de Comunicación, Riesgos y Documentación}}
\end{tabular}
\end{center}
\end{document}
